\section*{Overview}

Rollback DSU is the union-find structure for offline dynamic connectivity and other reversible merge problems. It keeps
union by size or rank, but deliberately gives up path compression so that every merge can be undone exactly.

\section*{When to Use It}

Use it when:

\begin{itemize}
  \item unions need to be undone later,
  \item the problem is offline,
  \item the solution explores a recursion tree over time intervals, query blocks, or divide-and-conquer branches.
\end{itemize}

\section*{Core Idea}

Every successful merge changes only a constant amount of state:

\begin{itemize}
  \item one parent pointer,
  \item one size or rank value,
  \item optionally one component counter or other component metadata.
\end{itemize}

Store those old values on a stack. A rollback simply pops and restores them.

\section*{Key Insight}

Path compression is the part that makes ordinary DSU hard to undo. It changes many parents at once. Rollback DSU avoids
that entire problem by keeping \texttt{find} reversible.

\section*{Worked Problem}

\paragraph{Problem.}
Edges of a graph are added and removed over time. For many time points, answer whether two vertices are connected.

\paragraph{Why rollback fits.}
The standard offline reduction places each edge on the segment tree over the time interval where that edge is active.
During the segment-tree DFS:

\begin{itemize}
  \item entering a node means applying all edges active on that whole interval,
  \item leaving the node means undoing those merges.
\end{itemize}

\section*{Correctness Intuition}

If you take a snapshot before entering one recursion branch and roll back when leaving it, then every branch sees the
exact DSU state corresponding to its own active edges. Because each merge records the precise state it overwrote,
rollback restores a valid earlier DSU state exactly.

\section*{Complexity Analysis}

\begin{itemize}
  \item \texttt{find}: \(O(\log n)\) amortized with union by size,
  \item \texttt{unite}: \(O(\log n)\),
  \item one rollback step: \(O(1)\),
  \item memory: \(O(n + \text{number of successful merges})\).
\end{itemize}

\section*{Implementation}

The code exposes:

\begin{itemize}
  \item \texttt{find},
  \item \texttt{unite},
  \item \texttt{same},
  \item \texttt{snapshot},
  \item \texttt{rollback}.
\end{itemize}

\section*{Common Pitfalls}

\begin{itemize}
  \item Adding path compression and destroying reversibility.
  \item Forgetting that failed merges still affect recursion logic if you do not record them consistently.
  \item Rolling back to the wrong snapshot boundary in a recursive solver.
  \item Treating rollback DSU as if it solved online edge deletions by itself.
\end{itemize}

\section*{Variants / Extensions}

\begin{itemize}
  \item Rollback DSU with component sums, parity, or bipartite flags.
  \item Segment tree over time for dynamic connectivity.
  \item Persistent DSU, which gives version access rather than explicit undo.
\end{itemize}

\section*{Practice Problems}

\begin{itemize}
  \item Offline dynamic connectivity.
  \item Bipartiteness under interval-active edges.
  \item Any divide-and-conquer over time where merges must be reversible.
\end{itemize}

\section*{References}

\begin{itemize}
  \item \href{https://cp-algorithms.com/data_structures/deleting_in_log_n.html}{cp-algorithms: Deleting from a Data Structure in \(O(T(n)\log n)\)}.
  \item \href{https://github.com/kth-competitive-programming/kactl}{KACTL}.
  \item \href{https://codeforces.com/catalog}{Codeforces Catalog}.
\end{itemize}
